\section{P. The full profinite comparison}

\begin{statement}{P.1}{Known}{What the opened comparison theorems assert}
Harbater--Schneps' Main Theorem identifies the symmetric peripheral outer
groups of the profinite fundamental groups of $\mathcal M_{0,n}$ with
$\widehat{\mathrm{GT}}_0$ for $n=4$ and with
$\widehat{\mathrm{GT}}$ for $n\geq5$.
It supplies an inclusion, without a strictness assertion.
Lochak--Schneps, Proposition 4, proves
$\widehat F_2^{c}=1$ and
$C_{\widehat{\mathrm{GT}}}(c)=\langle c\rangle$, where
$c=(-1,1)$. Its proof uses only the two cycle relations, so the
centralizer assertion also holds in $\widehat{\mathrm{GT}}_0$.
These assertions alone do not classify all involutions.
\end{statement}
\begin{proof}
The identification of the four-point pure mapping class group with
$\widehat F_2$ identifies its inertia classes with the three classes
in the question. The effective permutation action on them is $S_3$;
the kernel of the four-letter permutation action is the Klein group.
Thus the definition of the symmetric peripheral group in the Main
Theorem of \cite{PBC_HS2000} is precisely the question's definition.
The forgetful map from five to four points gives the stated inclusion.

For the second assertion the semidirect product
$\widehat F_2\rtimes\langle c\rangle$ is the free profinite product of
three groups of order two, generated by $c,xc,yc$. The finite-factor
centralizer theorem makes the centralizer of $c$ its own factor;
intersection with $\widehat F_2$ is trivial. This is the fixed-group
calculation in \cite[Proposition 4(i)]{PBC_LS1997}. If a normalized pair
$(\lambda,f)$ commutes with $c$, the multiplication of pairs gives
$f(x,y)=f(x^{-1},y^{-1})$, hence $f=1$. The proof of Proposition 4(ii)
then applies relation (II), and no pentagon, to obtain
$\lambda\in\{1,-1\}$. The same proof consequently applies to the gentle
group. Conjugacy of a different involution to $c$ is not a consequence
of a calculation of this centralizer.
\end{proof}

\begin{statement}{P.2}{Proved}{An integral non-pentagon special direction and its pro-prime lift}
For every prime $\ell\geq5$ there is an automorphism of the free
pro-$\ell$ group $F_\ell$ with normalized pair $(1,f_\ell)$ satisfying
the two cycle equations and failing the pentagon. We may prescribe
the degree-seven term of $f_\ell$ to be
\[
\begin{split}
 q={}&3b_{xxxyxyy}+3b_{xxxyyxy}+b_{xxyxyxy}\\
 &-3b_{xxyxyyy}-3b_{xxyyxyy}-3b_{xxyyyxy}-b_{xyxyxyy},
\end{split}
\]
where $b_w$ is the standard Lyndon bracket, with $x<y$.
This is a statement about actual continuous pro-$\ell$ automorphisms,
not just formal solutions of a linear equation.
\end{statement}
\begin{proof}
Write $L$ for the integral free Lie algebra on $x,y$ and put $z=-x-y$.
Our independent program \texttt{code/pbc\_q7.py} expands the displayed
brackets using $[a,b]=ab-ba$. It verifies, over $\mathbb Z$, the
identities
\[
 q(x,y)+q(y,x)=0,\qquad
 [x,q(x,y)]=[x+y,q(-x-y,y)],
\]
and the cyclic identity $q(x,y)+q(y,z)+q(z,x)=0$.
For clarity about the geometric content, set
$D(x)=0$, $D(y)=[y,q]$. The second identity gives
\[
 D(z)=[z,q(x,y)-q(z,y)].
\]
Thus $D$ is special. Under the transposition of $x,y$ the difference
from $D$ is the inner derivation $\operatorname{ad}(q)$; the same is
true under the cyclic permutation of $x,y,z$. The two displayed
identities verify these assertions on both generators. Its outer
class is therefore $S_3$-invariant.

Here is the integral lifting argument. For a special homogeneous
derivation of degree $n$, choose lifts in $\gamma_nF_\ell$ of the
three peripheral conjugators. The error in their product relation
lies in $\gamma_{n+2}F_\ell$. If the current error lies in
$\gamma_mF_\ell$, multiply the first two conjugators by elements of
$\gamma_{m-1}F_\ell$. The resulting change of the error in
$\gamma_m/\gamma_{m+1}$ is
\[
                  (u,v)\longmapsto[u,x]+[v,y].
\]
This is onto integrally: repeated use of the Jacobi identity expresses
every Lie monomial of degree $m$ as a sum of brackets with a generator.
We can cancel the error successively without changing the initial
derivation. The corrections converge in the pro-$\ell$ topology,
since every finite $\ell$-group is nilpotent. The resulting
endomorphism is the identity on the Frattini quotient, hence is onto
by the pro-$\ell$ basis theorem, and is an automorphism by finite
generation and the profinite Hopf property.

Let $P_n$ be the degree-$n$ filtration of the special outer
automorphism group. Its kernel on the Frattini quotient is pro-$\ell$:
the corresponding automorphism kernel of each finite $\ell$-group is
an $\ell$-group. An invariant element of $P_n/P_{n+1}$ has a lift
$a\in P_n$ whose discrepancies $a^{-1}g(a)$ form a cocycle in
$P_{n+1}$ for $g\in S_3$. This cocycle is a coboundary. Indeed, in
each finite stable quotient, complements to the normal $\ell$-group
are conjugate by the coprime complement theorem; equivalently one
solves successively along a central series, averaging over $S_3$.
The sets of corrections have the finite intersection property, and
compactness supplies a single correction in $P_{n+1}$. Here
$\ell\nmid6$ is essential. Consequently our direction has an actual
$S_3$-centralizing lift. Peripheral centralizers in a free pro-$\ell$
group give the unique normalization
$x\mapsto x$, $y\mapsto f_\ell^{-1}yf_\ell$, with
$f_\ell\in F_\ell'$. The generator calculations in
\cite[\S1.3, Propositions 3--4 and Theorem 2]{PBC_HS2000} use only
these centralizers and the free-generator substitutions. They give
exactly the two cycle equations in this completion as well.

We specify the pentagon certificate, including its comparison map.
Let $a,b,c$ freely generate $L_3$. For $i,j\leq3$ the derivation
$t_{ij}$ takes the $i$-th generator to its bracket with the $j$-th,
takes the $j$-th to its bracket with the $i$-th, and kills the third.
Put $t_{14}=\operatorname{ad}(a)$,
$t_{24}=\operatorname{ad}(b)$, $t_{34}=\operatorname{ad}(c)$.
All infinitesimal braid relations are verified in the program
(48 ordered tests). This is the conjugation representation of the
graded Lie algebra of the pure four-strand braid group on the graded
kernel of forgetting strand four. To obtain it directly, that kernel
is a free group of rank three, and the pure three-strand action is
trivial on its abelianization. The split extension is therefore an
almost direct product; induction on commutators gives its lower
central filtration and the displayed action. In particular a zero
graded braid element acts by the zero derivation.

With the defect convention
\[
\begin{split}
 \mathcal D(q)={}&q(t_{12},t_{23})
 +q(t_{12}+t_{13},t_{24}+t_{34})+q(t_{23},t_{34})\\
 &-q(t_{12},t_{23}+t_{24})-q(t_{13}+t_{23},t_{34}),
\end{split}
\]
the coefficient of $a^5bcb$ in $\mathcal D(q)(a)$ is $-16$.
Reversing the defect changes it to the $16$ of LM8 B.12.
There are $880,882,808$ nonzero coefficients in its three images,
respectively, for a total of $2570$. The full integer output is
saved. Thus this element is nonzero over $\mathbb Z_\ell$ for
every $\ell\geq5$. The first nonzero degree of the group pentagon
is its linear defect. Corrections in higher degrees cannot change
it, so the lifted element fails the pentagon.
\end{proof}

\begin{statement}{P.3}{Proved}{The central refinement argument supplies a full peripheral lift}
Let $\Lambda$ be either the free profinite group on two generators
with its three peripheral elements, or
$\widehat{C_2*C_3}$ with $s,r,(sr)^{-1}$.
In the second case let $U$ be the kernel of the standard quotient
to $S_3$; in the first case let $U=\Lambda$.
Every continuous automorphism of
$\Lambda/R_{\rm nil}(U)$ preserving these three positive
conjugacy classes lifts to a continuous peripheral automorphism of
$\Lambda$, with any three specified conjugators lifted.
\end{statement}
\begin{proof}
We give the finite argument used in LM7\_PUSH D.50, so that the
input supporting the next assertion is checked rather than imported
from the status of a previous return.

Take a finite quotient $q_E:\Lambda\twoheadrightarrow E$ and a
normal subgroup $A\triangleleft E$. Put $D=E/A$, $q_0=q_E\bmod A$,
$X_i=q_E(x_i)$, and let $T_i$ be the $A$-class of $X_i$.
For every subgroup $J\leq E$ mapping onto $D$ and meeting each
$T_i$, write $M=J\cap A$. For every $J$-invariant irreducible
complex representation $\rho$ of $M$, of dimension $d$, put
$e=\exp(J)$ and form the finite group
\[
 \widetilde J_\rho=\{(g,B):g\in J,
 B\rho(m)B^{-1}=\rho(gmg^{-1}),\ \det B\in\mu_e\}.
\]
Schur's lemma shows that its fibres over $J$ have exactly $de$
elements. The graph of $m\mapsto\rho(m)$ is a normal copy of
$M$; it meets the scalar kernel trivially. Its quotient $Q_\rho$
is a finite central extension of $D$ and
\[
                 \widetilde J_\rho=J\mathbin{\times}_D Q_\rho.
\]
For the free source, choose arbitrary lifts of the two generator
values to $Q_\rho$. For the modular source choose implementing
matrices with square, respectively cube, equal to the identity.
Such matrices exist by rescaling; their determinants belong to
$\mu_e$. Thus these choices define actual source homomorphisms
$R^0_\rho:\Lambda\to Q_\rho$, including in the modular case.
Their joint image with $q_0$ is a finite central extension
$D^\dagger\to D$, with an actual source map $q^\dagger$.

We check the comparison of proper lifting fibres. If $\chi$ is an
epimorphism from the free source to $D$, and
$V_i\subseteq J_{\chi(x_i)}$ are $M$-conjugacy invariant, their
triple-product count is
\[
 \frac1{|M|}
 \sum_{\rho\in\operatorname{Irr}(M)^J}
 \frac1{\dim\rho}
 \prod_{i=1}^3
 \sum_{a\in V_i}\operatorname{tr}\Psi_\rho(a,R_\rho(x_i)),
\]
where $\Psi_\rho$ is the matrix representation of
$J\times_DQ_\rho$ and $R_\rho$ is any source lift of $\chi$.
To prove the formula, put $B_i=\sum_{a\in V_i}a$ in
$\mathbb C[J]$. These sums commute with $\mathbb C[M]$.
In a primitive block, $B_i$ vanishes unless $\chi(x_i)$ fixes
that block; since the first two quotient values generate $D$,
only invariant representations survive. Each surviving matrix
sum is scalar by Schur's lemma. The identity coefficient of
$B_1B_2B_3$ is extracted by the regular character of $M$.
The three scalar denominators and the two dimension factors give
exactly the displayed formula. In the modular case the prescribed
classes already enforce the square and cube relations, so the same
count is a count of homomorphisms from the modular source.

Now take a prescribed epimorphism $q_1^\dagger:\Lambda\to
D^\dagger$ whose three values are conjugates of those of
$q^\dagger$, with specified conjugators $r_i$. Choose
$H_i\in E$ above their images in $D$. The new prescribed
$A$-classes are $T_i^1=H_iT_iH_i^{-1}$.
For each fixed $J$, choose $j_i\in J$ above the same base value.
Normality of $A$ gives the exact equality
\[
                 T_i^1\cap J=j_i(T_i\cap J)j_i^{-1}.
\]
The pair $(j_i,r_i\bmod Q_\rho)$ conjugates each old trace
sum to the new one. Thus the counts into every actual subgroup
$J$ agree, including subgroups with empty intersections, whose
counts are zero. M\"obius inversion on the finite subgroup
lattice now gives equality of the numbers of epimorphisms to
$E$. The old number is positive because of $q_E$. Hence there
is a new epimorphism to $E$ with the required three conjugators
modulo $A$. Notice that this part used only normality of $A$.

In the present application set
$A=\gamma_\infty(q_E(U))$. This subgroup is normal and satisfies
$[q_E(U),A]=A$; its quotient is nilpotent. The pure subgroup of
$D^\dagger$ is a central extension of this nilpotent group and
is therefore nilpotent. Hence every auxiliary source map factors
through $\Lambda/R_{\rm nil}(U)$. Precompose these maps with
the given automorphism of that quotient. We have precisely the
new epimorphism and conjugators required in the preceding count.

Finally impose, in the compact space $\Lambda^3$, the product
relation, generation in every finite quotient, and all prescribed
conjugator values in the retained quotient. Any finite collection
of these conditions is contained in one joint finite quotient
$E$, and the argument just proved meets it. Compactness supplies
three simultaneous conjugators satisfying every condition.
The source presentation gives an endomorphism; finite-quotient
generation makes it onto. A surjective endomorphism of a finitely
generated profinite group is injective: preimage permutes the
finitely many open subgroups of each bounded index, forcing a
kernel element into every open subgroup. This proves the assertion.
If a specified conjugator is $1$ in the retained quotient, an
inner correction in its kernel makes that conjugator literally $1$
without changing the retained automorphism.
\end{proof}

\begin{statement}{P.4}{Proved}{Strict containment for the full profinite groups}
There is a strict inclusion
\[
                    \widehat{\mathrm{GT}}
                    \varsubsetneqq\widehat{\mathrm{GT}}_0.
\]
In particular the conclusion concerns the full profinite groups,
and answers Task P.2 and the strictness corollary in Task C.5.
\end{statement}
\begin{proof}
Fix $\ell\geq5$ and take $(1,f_\ell)$ from P.2. Use the modular
convention $x=(sr)^{-2}$, $y=sxs$ for this proof. Relative to the
question this is $x_{\rm new}=y_{\rm old}^{-1}$ and
$y_{\rm new}=x_{\rm old}^{-1}$. Thus the change is the geometric
interchange composed with complex conjugation; it normalizes both
the gentle and pentagon groups and preserves satisfaction of the pentagon.
The two cycle relations define on
$B_\ell=\widehat B/O^\ell(K_2)$ the automorphism
\[
 s\longmapsto sf_\ell,\qquad
 r\longmapsto f_\ell^{-1}r,\qquad sr\longmapsto sr.
\]
Indeed its square and cube relations are the two cycles,
its restriction to the pure free pro-$\ell$ subgroup is the
given automorphism, and its map on $S_3$ is the identity.
Thus it is an automorphism of the extension as well. The two
torsion images have conjugators in the pro-$\ell$ kernel:
they are complements above the same subgroups of orders two
and three, so coprime complement conjugacy, followed by
compactness, supplies those conjugators.

The natural map
\[
 \widehat B/R_{\rm nil}(K_2)
       \longrightarrow\prod_{p,S_3}\widehat B/O^p(K_2)
\]
is an isomorphism. On the kernel this says that the maximal
pronilpotent quotient of $K_2$ is $\prod_pK_2^{(p)}$;
finite products are onto by coprime orders and Goursat's lemma,
and compactness gives the infinite product. Above a common
$S_3$-value, remove one lift and reduce to this kernel argument.
Use the preceding automorphism in the $\ell$-component and
the identity in all other components. The resulting actual
automorphism of the fibre product is peripheral and fixes $sr$.
Apply P.3 with the specified $sr$-conjugator equal to $1$.
Its normalized full lift $\Phi$ satisfies $\Phi(sr)=sr$.

Set $f=s\Phi(s)$. Its $S_3$-image and its pure abelianization
are trivial, since both are already retained by the pronilpotent
quotient. Hence $f\in K_2'$, and its $\ell$-image is exactly
$f_\ell$. The identities $\Phi(s)^2=\Phi(r)^3=1$ give
the two full cycle equations. Restriction to $K_2$ gives the
normalized pure map $x\mapsto x$, $y\mapsto f^{-1}yf$.
The action centralizes the outer $S_3$-action, since $\Phi$
is the identity on the quotient $S_3$. Thus $(1,f)$ belongs
to $\widehat{\mathrm{GT}}_0$.

If it belonged to $\widehat{\mathrm{GT}}$, its pentagon would
remain an identity after passage to the pro-$\ell$ braid
quotient. That identity has first nonzero graded defect
$\mathcal D(q)$, whose coefficient $-16$ was computed in P.2.
This contradiction proves strictness. The proof lifts the
entire pro-$\ell$ solution before using its defect; it does
not infer existence from a finite or graded residue alone.
\end{proof}

\begin{statement}{P.5}{Open}{The full literature inventory in Task P.1}
The requested inventory is not complete in this return.
The opened primary statements establish P.1 and the boundaries
recorded below; they do not compute the abelianization of
$\widehat{\mathrm{GT}}$, classify all of its finite characters,
or settle conjugacy of all its involutions.
\end{statement}
\begin{proof}[Verified comparisons and remaining step]
Nakamura--Schneps' paper explicitly introduces a subgroup by
(III$'$) and (IV), leaving its equality with the full group open
in that paper \cite[\S1]{PBC_NS2000}. The opened harmonic note
derives its conditions for Galois by comparing an open immersion
with a degree-two or degree-three map; it expressly leaves
properness of the resulting subgroup undecided
\cite[\S1]{PBC_NT2002}. Consequently neither source supplies
the substitution demanded in Task C.2 merely by naming the
orbifold map.

The supplied 2026 source of Lochak--Nakamura--Schneps, Theorem 1.1,
gives automorphism actions on the groups with zero or one marked
point. Its introduction expressly distinguishes such actions
from respecting the additional sewing rules. We read that
statement and this qualification; they cannot be substituted
for the missing genus-one comparison. The opened Schneps survey
asks the abelianization and involution-conjugacy questions,
and the contemporary sources checked here have not supplied
answers to them. Absence of an answer in this bounded search
is not a theorem that no later answer exists. The exact missing
work is the primary-text access gap recorded in P.7.
P.6--P.7 below complete the named Collas, Drinfeld, configuration,
and 2024--2026 mapping-class comparisons, and check Ihara's second
beta-function paper. No priority assertion is made for P.4
or for the supplied R24/R25.
\end{proof}

\input{sections/P_literature.tex}
